\section{Process-Level Evaluation}

\subsection{Process Evaluation Design}
\label{sec:process_evaluation_design}

Automated R\&D proceeds through repeated rounds in which an agent proposes a direction, implements the corresponding change, observes the result, and decides how to proceed. A final score alone cannot identify where this loop succeeds or fails.  We therefore decompose the process into Solution Framing (C1), Execution (C2), and Feedback Control (C3). This decomposition follows the causal structure of the loop: C1 evaluates what the agent chooses to pursue, C2 evaluates whether that choice is translated into a valid result, and C3 evaluates how subsequent decisions use experimental feedback. These stages represent distinct failure sources and therefore require different forms of improvement.

Unlike conventional one-shot tasks that provide feedback only on the final output, the iterative research tasks studied here expose explicit verifier feedback at each evaluated checkpoint. These step-level signals provide direct evidence of progress and failure, allowing us to evaluate the research process without relying on subjective model judgments.  We therefore compute all three scores deterministically from recorded evaluation signals, as detailed below. The resulting metrics are reproducible and auditable.

\textbf{C1: Solution Framing.}
C1 asks whether the directions an agent pursues lead quickly to a strong solution.
Rather than judging how sophisticated a proposal sounds, it uses the running best verifier score as an objective proxy for the quality of the directions discovered so far.
Trajectories are mapped to a common horizon, with shorter runs carrying their last running best forward and longer runs using a shared cutoff.
We summarize progress across the early, middle, and late parts of this horizon, so the score rewards both reaching a high score and reaching it early while preventing later failures from erasing an earlier discovery.

\textbf{C2: Execution.}
C2 asks whether an agent reliably translates proposed changes into executable and correct results.
At each non-initial evaluated checkpoint, a delivery gate first checks whether the artifact runs and, when the task provides a correctness verdict, whether it is correct.
Failed delivery receives no credit, while successful delivery is discounted according to the code-related build failures observed before that checkpoint.
The discount is bounded so that delivering a valid result remains the primary requirement, and failures caused by the environment are excluded.

\textbf{C3: Feedback Control.}
C3 asks whether an agent preserves successful discoveries and responds effectively when an attempted change makes the result worse.
Its retention component compares the final score with the highest step score reached during the run.
For each meaningful regression, its recovery component measures how much of the lost score is recovered and how many evaluated transitions the recovery requires, with additional self-evaluated attempts applying a bounded penalty for hidden trial and error.
The two components jointly reward preserving strong results and correcting setbacks; when no regression occurs, C3 uses retention alone because recovery was not tested.

We first average valid seeds within each pairing of model and task, and then weight tasks equally.  Appendix~\ref{app:process_metric_details} gives the complete formulas, hyperparameters, boundary rules, and reconstruction procedure.


\subsection{Process Capability Results}

\label{sec:process_results}

Using these three metrics, we compare capability profiles across models and task categories.

\begin{figure}[t]
\centering
\includegraphics[width=0.8\linewidth]{Figures/process_by_metric_axis.pdf}
\caption{Process dimensions across seven models. All values are averaged over three rollouts.}
\label{fig:process_dims}
\end{figure}

\textbf{Execution is broadly reliable, while Solution Framing and Feedback Control reveal greater variation.}
Opus-4.7 leads outcome at 0.739, C1 at 0.612, and C2 at 0.967, while also placing third on C3 at 0.920.  C2 is the most compressed dimension, ranging from 0.880 to 0.967, because successful delivery is common across the evaluated models.  C1 ranges from 0.473 to 0.612, while C3 ranges from 0.772 to 0.928.  The broader variation in C1 and C3 reveals differences that delivery success alone cannot explain.

\textbf{Similar outcomes can conceal sharply different Execution and Feedback Control profiles.}
GPT-5.5 and Gemini-3.1-Pro provide the clearest comparison between models with similar outcomes.  Their outcomes are 0.663 and 0.652, and both score 0.555 on C1, indicating nearly identical observed progress in solution framing. Their later capabilities differ sharply.  GPT-5.5 reaches 0.958 on C2 but 0.858 on C3, whereas Gemini-3.1-Pro reaches 0.889 on C2 but 0.920 on C3. Similar outcomes and framing quality can therefore arise from different balances between reliable implementation and feedback control. 
LongCat-2.0 provides a complementary perspective. Although it ranks sixth on outcome at 0.572 and on C1 at 0.478, it attains the highest observed C3 value at 0.928. This contrast shows that a lower overall outcome can conceal a relative strength in one part of the research process, reinforcing the value of examining process dimensions alongside final performance.

\noindent
\begin{minipage}[t]{0.56\linewidth}
\vspace{0pt}

\textbf{Different task categories expose different bottlenecks in the research loop.}
Figure~\ref{fig:process_by_cat} shows where each task category becomes constrained.  CUDA tasks has the lowest C1 at 0.370 and the lowest C2 at 0.850, but retains a high C3 of 0.924. Its main difficulty lies in discovering and implementing effective optimizations rather than preserving them once found. Model Development tasks shows the opposite pattern. It has the highest C2 at 0.985 but the lowest C3 at 0.743, indicating that runnable changes are easy to produce while optimization progress is harder to stabilize. Puzzle and Challenge tasks are strongest across the process, with C1 at 0.737 and both C2 and C3 near 0.930. These contrasts show that the same agent can face different bottlenecks depending on whether a task demands difficult solution discovery, reliable implementation, or stable response to feedback.

These headline scores identify where models and task categories differ. The trajectory diagnostics in the next section explain how those differences arise.
\end{minipage}\hfill
\begin{minipage}[t]{0.40\linewidth}
\vspace{0pt}
\centering
\includegraphics[width=\linewidth]{Figures/process_by_cat.pdf}
\captionsetup{type=figure,hypcap=false}
\caption{Process dimensions by task category, averaged over the seven models.}
\label{fig:process_by_cat}
\end{minipage}



\subsection{Behavioral Diagnostics}
\label{sec:process_further}

C1, C2, and C3 provide aggregate scores for three parts of the research loop. Figure~\ref{fig:process_diag} provides a more detailed view of the behaviors behind these scores, including how models make progress, implement changes, and respond to regressions. These diagnostics characterize behavior rather than form another overall ranking, so a larger value is not always better. The figure reports task balanced model averages, with the exact value printed in each cell. Color intensity indicates relative magnitude only within the same column, and the gray column reports the average number of evaluated commit rounds as observation support. Formal definitions and calculation details for all diagnostic measures are provided in Appendix~\ref{sec:appendix_diag}.

\begin{figure}[t]
\centering
\includegraphics[width=\linewidth]{Figures/process_diagnostics_compact.pdf}
\caption{Behavioral diagnostics across seven models. Each cell reports the exact value, while darker shading indicates a larger value within the same column and does not imply stronger capability. Ratios are shown as percentages, scores as decimals, and counts as averages. The gray column reports the average number of evaluated commit rounds as observation support. Values are averaged across repeated runs for each model and task.}
\label{fig:process_diag}
\end{figure}

\textbf{C1: routes to progress.}  
Best observed score reports the strongest evaluated solution reached during a run.  Early capture measures how much of that eventual peak is already present in the first evaluated round, while later headroom capture measures how much of the remaining score space is filled afterward.  Opus reaches the highest observed score at $0.757$ and records $53.4\%$ early capture and $53.0\%$ later headroom capture. Gemini-3.1-Pro reaches a lower best observed score of $0.667$, but combines the highest early capture at $83.7\%$ with the lowest later headroom capture at $16.5\%$.  GPT-5.5 begins at only $45.3\%$ of its eventual peak but later fills $46.9\%$ of the remaining score space. The three quantities distinguish the absolute quality of the best discovered solution, the strength of the initial direction, and subsequent progress.

\textbf{C2: implementation pathways.}  Builds per round measures the number of recognized build invocations observed before each evaluated round, while rounds with build errors reports the fraction of rounds containing at least one observed code-related build error.  Kimi-K2.7-Code and LongCat-2.0 obtain nearly identical C2 scores of $0.880$ and $0.888$, yet LongCat-2.0 performs $4.66$ builds per round, and encounters build errors in $17.1\%$ of rounds, compared with $2.70$ and $8.5\%$ for Kimi-K2.7-Code.  Similar delivery reliability can therefore conceal substantially different amounts of observable construction and repair.  GPT-5.5 and Gemini-3.1-Pro occupy the two extremes.  GPT-5.5 records only $0.51$ builds per round and build errors in $0.8\%$ of rounds, whereas Gemini-3.1-Pro records $7.49$ and $17.6\%$ while achieving a lower C2 score.  Dense build and repair activity before commit therefore does not by itself imply reliable delivery.  Claude-Opus-4.7 provides a more balanced reference, combining $2.17$ builds per round and build errors in only $3.9\%$ of rounds with the highest C2 score.  These diagnostics describe visible implementation pathways rather than the quality of the underlying reasoning.

\textbf{C3: feedback behavior and exposure.}  Peak retention measures how much of the best observed score is preserved in the final result. Dip rate and dip depth describe the frequency and severity of regressions, while recovery credit measures how completely and quickly the agent recovers, including a bounded penalty for additional evaluated candidates between commit rounds. Opus-4.7 and GLM-5.2 show the most balanced profiles.  Their peak retention values are $0.981$ and $0.958$, and their recovery credit values are $0.711$ and $0.703$, while both experience relatively shallow dips.  GPT-5.5 retains $0.959$ of its peak but has the highest dip rate at $0.134$.  Across an average of $10.12$ evaluated commit rounds, it experiences regressions more frequently but still obtains $0.614$ recovery credit. Gemini-3.1-Pro and LongCat-2.0 retain $0.988$ and $0.962$ of their peaks and record lower dip rates of $0.069$ and $0.052$, but their recovery credit values are only $0.323$ and $0.520$. They average just $2.54$ and $5.42$ evaluated commit rounds, so their low dip frequencies must be interpreted with their more limited exposure to regression. Their high C3 scores therefore arise mainly from peak retention and fewer observed regressions rather than a well supported recovery advantage. DeepSeek-V4-Pro has the lowest peak retention and the deepest dips, showing that its feedback control is limited by both loss of strong intermediate results and more severe regressions.  Evaluated commit rounds are reported as evidence rather than as an additional capability measure.

Together, the retained diagnostics answer distinct questions about discovered solution quality, initial direction, subsequent gains, implementation activity, retention, regression, and repair. They explain the process scores while keeping observation support separate from the scored dimensions.
